\subsection{Normal Sampling Pivots}
% Source integration: supplied Chapter 6 screenshots, Sections 6.2--6.4.

For normal samples, the sample mean and sample variance separate exactly.  This
is the finite-sample source of the classical $z$, $\chi^2$, and Student $t$
pivots used in normal-theory confidence intervals and tests.

\begin{definition} \label{def:normal-sample-mean-variance} \hlt{Normal Sample Statistics}
\begin{enumerate}[label=(\roman*)]
\setlength{\itemsep}{0pt}
\item \hlt{Sample}: $X_1,\ldots,X_n$ are independent with $X_i\sim N(\mu,\sigma^2)$, $n\geq2$, and $\sigma^2>0$.
\item \hlt{Sample mean}: $\overline{X}\defeq n^{-1}\sum_{i=1}^nX_i$.
\item \hlt{Sample variance}: $S^2\defeq(n-1)^{-1}\sum_{i=1}^n(X_i-\overline{X})^2$.
\end{enumerate}
\end{definition}

\begin{proposition} \label{prop:normal-sample-mean-law} \hlt{Normal Sample Mean}\\
Let $X_1,\ldots,X_n$ be independent with $X_i\sim N(\mu,\sigma^2)$, and set $\overline{X}\defeq n^{-1}\sum_{i=1}^nX_i$. Then $\overline{X}\sim N(\mu,\sigma^2/n)$.
\begin{proof} Since independent normal sums are normal, $\overline{X}$ is normal.\\
Also $\mathbb{E}\overline{X}=n^{-1}\sum_i\mathbb{E}X_i=\mu$ and $\operatorname{Var}(\overline{X})=n^{-2}\sum_i\operatorname{Var}(X_i)=\sigma^2/n$. Thus $\overline{X}\sim N(\mu,\sigma^2/n)$.\\
\hlt{Structural Commentary}: Averaging preserves normality and divides variance by $n$.\\
\hlt{Key Insight}: The sample mean is an exact normal pivot center, not an asymptotic approximation.
\end{proof}
\end{proposition}

\begin{theorem} \label{thm:normal-mean-residual-independence} \hlt{Mean-Residual Independence}\\
Let $X_1,\ldots,X_n$ be independent with $X_i\sim N(\mu,\sigma^2)$, set $\overline{X}\defeq n^{-1}\sum_iX_i$, $R\defeq(X_1-\overline{X},\ldots,X_n-\overline{X})$,\\
and $S^2\defeq(n-1)^{-1}\sum_iR_i^2$. Then $\overline{X}\perp\!\!\!\perp R$, hence $\overline{X}\perp\!\!\!\perp S^2$.
\begin{proof} Let $Z_i\defeq(X_i-\mu)/\sigma$, so $Z\sim N_n(0,I_n)$.\\
Choose orthogonal $Q$ with first row $n^{-1/2}\mathbf{1}^\top$, set $Y\defeq QZ$. Then $Y\sim N_n(0,I_n)$, $Y_1\perp\!\!\!\perp(Y_2,\ldots,Y_n)$.\\
Since $\overline{X}=\mu+\sigma Y_1/\sqrt n$ and $R=\sigma(I-n^{-1}\mathbf{1}\mathbf{1}^\top)Z=\sigma Q^\top(0,Y_2,\ldots,Y_n)^\top$, $\overline{X}$ is a function of $Y_1$ and $R$ is a function of $(Y_2,\ldots,Y_n)$. Therefore $\overline{X}\perp\!\!\!\perp R$, and $S^2=(n-1)^{-1}\sum_iR_i^2$ gives $\overline{X}\perp\!\!\!\perp S^2$.\\
\hlt{Structural Commentary}: Mean and residuals are orthogonal projections.\\
\hlt{Key Insight}: For Gaussian vectors, orthogonal components are independent.
\end{proof}
\end{theorem}

\begin{theorem} \label{thm:normal-sample-variance-chi-square} \hlt{Sample Variance Chi-Square Law}\\
Given independent $X_1,\ldots,X_n$; $X_i\sim N(\mu,\sigma^2)$;$S^2\defeq(n-1)^{-1}\sum_{i=1}^n(X_i-\overline{X})^2$. Then $(n-1)S^2/\sigma^2\sim\chi^2_{n-1}$.
\begin{proof} Let $Z_i\defeq(X_i-\mu)/\sigma$, so $Z=(Z_1,\ldots,Z_n)\sim N_n(0,I_n)$. Choose an orthogonal matrix $Q$ whose first row is $n^{-1/2}(1,\ldots,1)$, and set $Y\defeq QZ$; orthogonal invariance gives $Y\sim N_n(0,I_n)$ with independent coordinates.\\
Then $Y_1=\sqrt n(\overline{X}-\mu)/\sigma$ and $\sum_{j=2}^nY_j^2=\sum_{i=1}^n(Z_i-\overline{Z})^2=(n-1)S^2/\sigma^2$.\\
Thus $(n-1)S^2/\sigma^2$ is the sum of squares of $n-1$ independent $N(0,1)$ variables, hence $\chi^2_{n-1}$.\\
\hlt{Structural Commentary}: Centering removes the one-dimensional mean direction and leaves $n-1$ residual directions.
\hlt{Key Insight}: The sample variance is the squared Gaussian residual length scaled by $\sigma^2$.
\end{proof}
\end{theorem}

\begin{corollary} \label{cor:normal-student-t-pivot} \hlt{Student Pivot for a Normal Mean}\\
Let $X_1,\ldots,X_n$ be independent with $X_i\sim N(\mu,\sigma^2)$, and set $\overline{X}\defeq n^{-1}\sum_iX_i$ and $S^2\defeq(n-1)^{-1}\sum_i(X_i-\overline{X})^2$.\\
Then $(\overline{X}-\mu)/(S/\sqrt n)\sim t_{n-1}$.
\begin{proof} Let $Z\defeq\sqrt n(\overline{X}-\mu)/\sigma$ and $V\defeq(n-1)S^2/\sigma^2$.\\
By normal mean law, $Z\sim N(0,1)$. By sample-variance chi-square law, $V\sim\chi^2_{n-1}$.\\
By mean-residual independence, $Z\perp\!\!\!\perp V$. Therefore $(\overline{X}-\mu)/(S/\sqrt n)=Z/\sqrt{V/(n-1)}\sim t_{n-1}$.\\
\hlt{Structural Commentary}: Replacing $\sigma$ by the independent estimate $S$ changes a normal pivot into a Student pivot.
\hlt{Key Insight}: Student's $t$ is a normal variable divided by an independent chi-square scale estimate.
\end{proof}
\end{corollary}

\begin{method} \label{meth:normal-mean-exact-confidence-intervals} \hlt{Exact Normal Mean Intervals}\\
Let $X_1,\ldots,X_n$ be independent $N(\mu,\sigma^2)$ variables, with $\overline{X}\defeq n^{-1}\sum_iX_i$ and $S^2\defeq(n-1)^{-1}\sum_i(X_i-\overline{X})^2$.\\
Below cases have exact coverage $1-\alpha$.
\begin{enumerate}[label=(\roman*)]
\setlength{\itemsep}{0pt}
\item \hlt{Known variance}: if $\sigma$ is known, $C_\mu(X)=[\overline{X}-z_{1-\alpha/2}\sigma/\sqrt n,\overline{X}+z_{1-\alpha/2}\sigma/\sqrt n]$
\item \hlt{Unknown variance}: if $\sigma$ is unknown, $C_\mu(X)=[\overline{X}-t_{n-1,1-\alpha/2}S/\sqrt n,\overline{X}+t_{n-1,1-\alpha/2}S/\sqrt n]$
\item \hlt{Variance}: $C_{\sigma^2}(X)=[(n-1)S^2/\chi^2_{n-1,1-\alpha/2},(n-1)S^2/\chi^2_{n-1,\alpha/2}]$
\end{enumerate}
\end{method}
